\documentclass[11pt,a4paper,openany]{book}

\makeatletter\def\input@path{{../_shared/}}\makeatother
\usepackage{course}

\begin{document}
\raggedbottom

\coursetitle{Game Theory and Strategic Decisions}{BA5}{2026--2027}

\tableofcontents
\clearpage

These notes cover the course by Elena Perazzi (EPFL, Fall 2026): how to model strategic
interactions as games, and how to predict their outcome with solution concepts --- Nash
equilibrium, dominance, subgame-perfect and Bayesian Nash equilibria.


% =========================================================
\chapter{Lecture 1 : Games and Nash equilibrium}
% =========================================================

\begin{rembox}[title=In a nutshell]
\textbf{What is this lecture about?} What a game is (players, choices, payouts), how to write
a simultaneous game as a table (\emph{normal form}), the difference between actions and
(pure or mixed) strategies, and the central solution concept of the course: the
\textbf{Nash equilibrium}. \\
\textbf{Why is it useful?} Every later lecture refines or reuses these notions. The classical
games seen here (Prisoner Dilemma, Matching Pennies, Split or Steal) are the reference
examples of the whole course.
\end{rembox}

\begin{rembox}[title=Course organisation]
\begin{itemize}
  \item 14 lectures, on Tuesday from 1:15pm to 3pm (one 15-minute break around 2pm).
  \item Assistant: Luiz Bissoto (\href{mailto:luiz.bissoto@epfl.ch}{luiz.bissoto@epfl.ch}).
  \item Grading based on \textbf{3 problem sets}, weights \textbf{40\%, 40\%, 20\%};
        assignments can be done in groups of up to 3 people.
  \item Three parts: basics of non-cooperative game theory (approx.\ 7 lectures), auction
        theory with applications to Ebay and Google advertising auctions (approx.\ 4
        lectures), evolutionary game theory (approx.\ 2 lectures).
\end{itemize}
\end{rembox}


\section{Cooperative vs non-cooperative games}

\begin{defnbox}[title=Non-cooperative game]
A game in which there is \textbf{no external authority enforcing rules of cooperative
behavior}. Each player tries to maximize her benefit at any time, even by damaging others.
Even when a contract between two players could be beneficial to both, they cannot stipulate
it, because no authority is able to enforce it.
\end{defnbox}

\begin{defnbox}[title=Cooperative game]
A game that involves negotiation, bargaining, contracts, \ldots
\end{defnbox}

\begin{rembox}[title=Scope of the course]
Classical game theory focuses on \textbf{non-cooperative games}: this is what we study in
this course.
\end{rembox}

\end{document}
